\documentclass[aps,prb,amsmath,amssymb]{revtex4}
\usepackage{bm}
\usepackage{graphicx}

\DeclareMathOperator{\sgn}{sgn}
\DeclareMathOperator{\Img}{\mathrm{Im}}
\DeclareMathOperator{\Rea}{\mathrm{Re}}
\DeclareMathOperator{\Sp}{\mathrm{Sp}}
\DeclareMathOperator{\tr}{\mathrm{tr}}
\DeclareMathOperator{\arctanh}{arctanh}



\begin{document}


\title{Nonequilibrium Josephson Effect in SIS Junctions}
\author{Mikhail S. Kalenkov}
\affiliation{I.E. Tamm Department of Theoretical Physics, P.N. Lebedev Physical Institute, 119991 Moscow, Russia}

\date{\today}
\begin{abstract}

\end{abstract}
\maketitle

This document summarizes the equations solved numerically in the code and specifies the conventions used throughout the calculations.

\section{Model}

The nonequilibrium Josephson current in an SIS junction can be represented in a compact matrix form (for details see Ref.~\onlinecite{KZ24})
\begin{equation}
I(t) = 
\dfrac{e}{8}
\Sp \Bigl\{
[\mathbb{S}_n + \Gamma]^{-1}
(H - \Gamma H \Gamma^+)
\\\times
[\mathbb{S}_n^+ + \Gamma^+]^{-1}(\mathbb{S}_n^+
\hat \tau_3 \hat o_3
\mathbb{S}_n - \hat \tau_3 \hat o_3)
\Bigr\},
\end{equation}
where $I$ denotes the Josephson current carried by a single conducting channel that hosts quasiparticle states with both spin orientations, $\hat \tau_3$ and $\hat o_3$ are Pauli matrices in Nambu and trajectory spaces, respectively, and $\mathbb{S}_n$ is the normal-state interface scattering matrix.

For simplicity we assume that the interface preserves the quasiparticle spin projection. Then the scattering matrix for a given spin projection $\sigma=\pm1$ has the form
\begin{equation}
[\mathbb{S}_n]_{\sigma}
=
\begin{pmatrix}
\sqrt{R_{\sigma}}  & -i \sqrt{D_{\sigma}}  &  0 & 0
\\
-i \sqrt{D_{\sigma}}  & \sqrt{R_{\sigma}}  & 0 & 0 \\
0 & 0 & \sqrt{R_{-\sigma}}  & i \sqrt{D_{-\sigma}}
\\
0 & 0 & i \sqrt{D_{-\sigma}}  & \sqrt{R_{-\sigma}}  \\
\end{pmatrix}
e^{-i\sigma\theta/2},
\end{equation}
is the normal-state interface transmission probability for quasiparticles with spin projection $\sigma$, and $\theta$ is the spin-mixing angle.

For a given time-dependent voltage bias $V(t)$, the superconducting phase difference across the junction reads
\begin{equation}
\varphi (t)= 2e\int^tV(t')dt'.
\label{phi}
\end{equation}
The kernels of the operators $\Gamma$ and $H$ are then given by
\begin{gather}
\Gamma(t,t') = e^{i\varphi(t) \hat \tau_3 \hat o_3/4} \Gamma_{eq}(t-t') e^{-i\varphi(t') \hat \tau_3 \hat o_3/4},
\\
H(t,t') = e^{i\varphi(t) \hat \tau_3 \hat o_3/4} H_{eq}(t-t') e^{-i\varphi(t') \hat \tau_3 \hat o_3/4},
\end{gather}
where 
\begin{gather}
H_{eq}(t-t')
=
\int_{-\infty}^{\infty}\dfrac{d \varepsilon}{2\pi}
h_0(\varepsilon) e^{-i\varepsilon(t - t')},
\quad h_0(\varepsilon) = \tanh\dfrac{\varepsilon}{2\pi},
\\
\Gamma_{eq}(t-t')
=
\hat \tau_1
\int_{-\infty}^{\infty}\dfrac{d \varepsilon}{2\pi}
a(\varepsilon) e^{-i\varepsilon(t - t')},
\quad a(\varepsilon) = \dfrac{-\varepsilon + \sqrt{\varepsilon^2 - |\Delta|^2}}{|\Delta|}.
\end{gather}
We assume that the superconducting order parameter is spatially uniform in both leads and has the same magnitude $|\Delta|$, corresponding to a symmetric junction.

For a time-independent voltage bias, evaluation of the current can be reduced to solving matrix recurrence relations derived in Ref.~\onlinecite{KZ24}. These recurrence relations are implemented numerically in the present code.

For reference, in the limit of a transparent junction biased by a small voltage $|eV| \ll |\Delta|$ one obtain \cite{AB1} 
\begin{equation}
I = e |\Delta|  |\sin(eVt)| \tanh \dfrac{|\Delta|}{2T} \sgn V.
\end{equation}




\begin{thebibliography}{99}
\bibitem{KZ24} M.S. Kalenkov and A.D. Zaikin, Absence of fractional ac Josephson effect in superconducting junctions,
Phys. Rev. B 110, 134521 (2024).
\bibitem{AB1} D. Averin and A. Bardas, ac Josephson Effect in a Single Quantum Channel, Phys. Rev. Lett. \textbf{75}, 1831 (1995).
\end{thebibliography}

\end{document}